\section*{Limitations}
Our study faced several significant limitations that impacted our approach and findings. A primary constraint was the token limitation and high subscription charges for paid cloud services, which restricted our ability to perform inference and fine-tuning on larger parametric models, particularly those with 70B or 40B parameters. This limitation likely curtailed our exploration of the full capabilities of these advanced models, which could have provided deeper insights or enhanced performance.

Another critical limitation was the resource-intensive nature of obtaining legal expert annotations. Due to the high costs and extensive time required for this process, it was not feasible for us to obtain expert evaluations for the entire PredEx test dataset. Consequently, we opted to sample 50 random documents for expert review and Likert score evaluations. While necessary, this approach potentially limits the breadth and depth of our expert-based evaluation, as it does not encompass the entire dataset.

In terms of the effectiveness of Large Language Models (LLMs) in the legal domain, our findings suggest that while these models are proficient in conversational contexts, their applicability in logic or knowledge-intensive tasks like legal judgment prediction and explanation is less convincing. Analyzing lengthy legal documents and generating predictions with explanations poses a significant challenge for generative-based models. This is particularly true in cases where the models need to process and understand complex legal reasoning and contexts.

Furthermore, the performance of the open-source baseline model, which was intended to jointly predict and generate explanations, did not meet our expectations. This underperformance could be attributed to the token limitations imposed during our study. By only using the last 1000 tokens of documents for fine-tuning, there is a possibility that the model did not fully grasp the entire context of the cases. Moreover, our fine-tuned models frequently produced truncated responses due to the 512-token limit set for generation. This limitation may have hindered the models' ability to generate comprehensive and nuanced explanations.

Lastly, the pre-trained models used in our study inherently lacked detailed knowledge specific to Indian legal cases. Even after undergoing tuning processes, these models struggled to generate explanations that paralleled the depth and specificity of human-like legal reasoning. This shortfall highlights the challenge of adapting general AI models to specialized domains such as law, where domain-specific knowledge and reasoning are crucial.

These limitations underscore the challenges in applying LLMs to complex and specialized tasks like legal judgment prediction and explanation. They also highlight the necessity for continued research and development efforts aimed at enhancing the capabilities of AI models in interpreting and understanding legal documents and contexts.
